%%%%%%%%%%%%%%%%%%%%%%%%%%%%%%%%%%%%%%
% Jordan Moak's Resume               %
% Type: Software Engineer            %
%%%%%%%%%%%%%%%%%%%%%%%%%%%%%%%%%%%%%%

%!TEX TS-program = xelatex
\documentclass[]{moak-resume}
% All packages are loaded by moak-resume.cls

% Document Meta information.
\hypersetup{
    pdftitle={Jordan Moak - Software Engineer, Platform and AI},
    pdfauthor={Jordan Moak},
    pdfsubject={Resume of Jordan Moak, Software Engineer specializing in Platform Engineering and AI},
    pdfkeywords={Software Engineer, Platform Engineering, Developer Tooling, AI Engineering, Release Engineering, Deployment Automation, CI/CD, Jenkins, Backstage, SonarQube, Akamai, Entra ID, Kubernetes, AWS, MCP, Model Context Protocol, AI Enablement, Agentic Development, Coding Agents, Context Engineering, Cursor, opencode, LangGraph, LangChain, Local LLMs, Agent Memory, OpenClaw, TypeScript, Python, Go, Java, SQL, Bash, Groovy, PowerShell, SAS, Node.js, Bun, Event-Driven Architecture, Microservices, Test Automation, Appium, Selenium},
    colorlinks=false,           % no link border color
    allbordercolors=white,      % white border color for all
    pdfcreator={XeLaTeX}
}

\begin{document}

\makeprofile{Jordan}{Moak}
            {Software Engineer, Platform \& AI}
            {\icon{\faPhone} +513 417 4499}
            {\icon{\faEnvelope[regular]} \href{mailto:jordan.moak@gmail.com}{ jordan.moak@gmail.com}}
            {\icon{\faGlobe} \href{https://moak.dev}{ \contactprimary{moak.dev}}}
            {\icon{\faGithub} \href{https://github.com/JMoak/}{ JMoak}}
            {\icon{\faLinkedin} \href{https://www.linkedin.com/in/jmoak/}{ /in/JMoak}}
            {\icon{\faMapMarker} Loveland, OH}

% In the aside, each new line forces a line break
\begin{aside}
  \profilesection{About Me}
  \small\bodyfont Since 2013 I have built the things that let other people ship: deployment pipelines, test frameworks, release tooling. Now I build them for coding agents: MCP servers, and keywork, a workspace for running agents. I believe a change should leave the next one easier to make.
  \profilesection{Focus}
	\focusChart{}
  ~
  \profilesection{Technical Skills}
  \skillcategory{01}{Languages}{TypeScript \skillsep Python \skillsep Go \skillsep Java\\SQL \skillsep Bash \skillsep Groovy\\PowerShell \skillsep SAS}
  \skillcategory{02}{Platform}{Release Automation \skillsep CI/CD\\Jenkins \skillsep Backstage \skillsep SonarQube\\Akamai \skillsep AWS \skillsep K8s}
  \skillcategory{03}{AI Engineering}{MCP Servers \skillsep Coding Agents\\Context Engineering \skillsep LangGraph\\Local LLMs \skillsep Agent Memory}
\end{aside}
\maincontentstart
\section{Exp|erience}

  % -- Care.com
  \begin{entrylist}
    \experienceentry
    {[Oct 2021 - Present]}
    {Senior Software Engineer}
    {Care.com}
    {}
  \end{entrylist}
  \begin{experienceitems}
    \item Automated the critical monolith's deployment pipeline in Groovy/Jenkins across engineering: biweekly 140+ commit batches became twice-weekly releases of 10--70 commits in late 2023, then weekday deployments from 2025, with no engineer assigned to releases.
    \item Production P0 and P1 incidents fell 90\% from 2023 to 2024, largely through fewer human errors in releases, returning engineering time to development.
    \item Developed event-driven systems processing billions of consumer-data records annually, supporting CRM campaigns and lifecycle messaging.
    \item Lead the company's AI enablement since early 2025: introduced Cursor and built context to help coding agents navigate and work in the critical monoliths.
    \item Lead MCP development: built the first internal MCP servers and enabled adoption across engineering through reusable patterns and hands-on guidance.
    \item Built a domain-expert agent prototype and a prototype MCP connector for flagship AI assistants.
    \item Led a proof of concept for a containerized coding agent on opencode meeting company security standards; presented model cost comparisons to senior leadership.
    \item Unblocked the Backstage service catalog with Active Directory and Entra ID sign-in and SonarQube in CI/CD; established platform-agnostic patterns adopted as organization-wide standards.
    \item Completed a targeted replatform from Apache to Akamai in two weeks, saving \$48K+ annually.
    \item Restored critical features and CI/CD for the legacy platform supporting a 10-person senior engineering team.
    \item Gave recurring department-wide talks on platform changes, engineering methods, and rollouts.
  \end{experienceitems}
  \experiencegap

 % -- ANGI
 \begin{entrylist}
  \experienceentry
    {[Nov 2018 - Oct 2021]}
    {Software Engineer III}
    {Angi (formerly known as Angie's List)}
    {}
\end{entrylist}
\begin{experienceitems}
  \item Built code-generation templates and developer tooling adopted by delivery teams across Angi.
  \item Designed a testing framework for REST microservices on Node.js, Mountebank, and Cucumber, with CI/CD on CircleCI.
  \item Led development of an Appium mobile testing framework in Python, with XCUITest and UIAutomator2, for four mobile apps.
  \item Built and maintained a Selenium frontend testing framework for a microfrontend architecture.
  \item Taught automation practices through bootcamps, documentation, and organizational support.
\end{experienceitems}
\experiencegap

 % -- ROCHE
\begin{entrylist}
  \experienceentry
    {[Sep 2014 - Nov 2018]}
    {DV Software Engineer}
    {Roche Diagnostics (Contracted through PRO Unlimited)}
    {}
\end{entrylist}
\begin{experienceitems}
  \item Diagnosed and fixed a floating-point precision defect, restoring valid results and returning a production line with millions of dollars in daily output to operation.
  \item Automated FDA-compliant validation of a mobile, web, and medical-device system in Java, Appium, Selenium, and Bash on Linux.
  \item Built reusable automation tools in Java, Selenium, PowerShell, and SAS within regulated SOPs.
  \item Led team automation of test datasets and configuration files for medical software verification.
\end{experienceitems}
\experiencegap

% -- TSC
\begin{entrylist}
  \experienceentry
    {[Mar 2014 - Aug 2014]}
    {Programmer/Analyst}
    {The Stevenson Company}
    {}
\end{entrylist}
\begin{experienceitems}
  \item Wrote statistical procedures for quality analysis and data cleaning in market research web apps.
  \item Migrated SAS programs to MySQL and PHP.
\end{experienceitems}
\experiencegap

% -- STATKING
\begin{entrylist}
    \experienceentry
    {[Jan 2013 - Dec 2013]}
    {Statistical Programmer}
    {STATKING Clinical Services}
    {}
  \end{entrylist}
\begin{experienceitems}
  \item Led SAS programming for 20+ Phase I-III clinical trials, developing CDISC SDTM/ADaM datasets.
  \item Produced 160+ analysis tables and listings for a Phase III clinical trial submission.
\end{experienceitems}

\makecontinuationpage{Jordan}{Moak}{Software Engineer, Platform \& AI}

\section{Sel|ected Personal Projects}
\begin{entrylist}
  \entry
    {[GitHub: JMoak]}
    {\href{https://github.com/JMoak/keywork}{keywork} and \href{https://github.com/JMoak/keywork-app}{keywork-app}}
    {Terminal and Desktop Agent Workspaces}
    {TypeScript, Bun, Electron, Solid; public releases}
\end{entrylist}
\begin{itemize}
  \item Built keyboard-first terminal and desktop workspaces for coding agents, with tiled sessions, branching conversations, and a permission gate that requires approval for edits and commands by default.
  \item A shared headless engine serves the terminal, CI runs, and desktop through one typed event stream. File-editing turns are snapshotted in git for undo; desktop approval cards show diffs, commands, and tool arguments before execution.
  \item Released as single-file binaries for Linux, macOS, and Windows and as keywork-cli on npm; keywork-app ships installers for all three.
\end{itemize}

\begin{entrylist}
  \entry
    {[npm: @jmoak]}
    {\href{https://github.com/JMoak/ratchet}{ratchet}}
    {Decision Engine for Model Judgments}
    {}
\end{entrylist}
\begin{itemize}
  \item Built a TypeScript decision library that pools model judgments with a correction for correlated errors; deterministic rules can require review or refusal regardless of model confidence.
  \item Decisions are logged in shadow mode before a gate may act; promotion requires passing a predefined statistical test on labeled outcomes. Published on npm with an experimental API; operational accuracy not yet established.
\end{itemize}

\begin{entrylist}
  \entry
    {[npm \& PyPI]}
    {MCP and Testing Packages}
    {Developer Tooling}
    {}
\end{entrylist}
\begin{itemize}
  \item \textbf{chrono-mcp} (@jmoak/chrono-mcp): an MCP server for current time, timezone conversion, and batch date arithmetic.
  \item \textbf{vibrissa} (@jmoak/vibrissa): records real MCP server responses as JSON cases and replays them in CI; chrono-mcp's CI uses it.
  \item \textbf{tdd-dsl}: a test-driven development DSL, published to npm and PyPI.
\end{itemize}

\begin{entrylist}
  \entry
    {[GitHub: Rosavera-I]}
    {\href{https://github.com/Rosavera-I}{rosavera}}
    {OpenClaw Personal Assistant}
    {}
\end{entrylist}
\begin{itemize}
  \item An OpenClaw personal assistant in daily use for maintaining code projects and trying agent ideas.
  \item Federated memory: the people in a conversation decide which vaults retrieval may search and where its transcript is stored; embeddings run locally through Ollama, indexed in SQLite with sqlite-vec.
\end{itemize}

\begin{entrylist}
  \entry
    {[GitHub: JMoak]}
    {\href{https://github.com/JMoak/nrly}{nrly}}
    {Weight-Editing Workbench}
    {Experimental; source only, no release}
\end{entrylist}
\begin{itemize}
  \item Built a local workbench for reversible, reproducible model-weight experiments: a Python/PyTorch kernel, CLI, 26-tool MCP server, and typed TypeScript client.
  \item Recorded and replayed a Qwen3-0.6B weight-editing study with exact restoration of the original model; results were inconclusive against a simple prompt-change control.
\end{itemize}

\begin{entrylist}
  \entry
    {[Private repo]}
    {\href{https://moak.dev/stead}{stead}}
    {Personal Agent Desktop}
    {In development; private repository, no release}
\end{entrylist}
\begin{itemize}
  \item Built a separate Electron app on keywork-app's SDK and keywork's runtime, presenting personal agents as creatures in a homestead.
  \item Designed a visual language for agent activity, pending approvals, and returned work, with scoped permissions that keep authorization with the user.
\end{itemize}

\section{Edu|cation}
\begin{entrylist}
  \entry
    {[2008 - 2013]}
    {Bachelor's Degree in Mathematics}
    {Northern Kentucky University}
    {Specialization in Applied Mathematics}
  \entry
    {[2008 - 2013]}
    {Bachelor's Degree in Statistics}
    {Northern Kentucky University}
    {Capstone: R-based anomaly detection for credit card fraud}
\end{entrylist}
\end{document}
